% TOP_PROBLEM: {"id":30007746,"problem_number":"LOCAL-30007746","title":"Scaling housing mobility programs","category_id":21,"category":{"id":21,"name":"economics","display_name":"Economics","description":"Open economic theory statements, published research agendas, and proposed scoped specifications; question provenance and historical priority are qualified per record.","slug":"economics","order_index":21,"created_at":"2026-10-08T00:00:00Z"},"status":"research_question","external_url":null,"legacy_ids":[],"published":false,"statement":"Estimate welfare at different saturation levels of a housing-mobility assistance program after rents, neighborhood composition, and school capacity respond.","economics":{"original_id":235,"field":"Urban, housing and spatial economics","kind":"empirical","formalization_basis":"adapted_research_question","source_status":"explicit_open","priority_status":"earliest_found","priority_note":"Earliest explicit statement located in this audit; no exhaustive priority search or first-ever claim. The mathematical specification below is adapted, not quoted from the source.","earliest_verified_reference":{"authors":["Eric Chyn","Lawrence F. Katz"],"title":"Neighborhoods Matter: Assessing the Evidence for Place Effects","year":2021,"url":"https://lkatz.scholars.harvard.edu/sites/g/files/omnuum5961/files/lkatz/files/chyn_katz_jep_2021_all.pdf","doi":"10.1257/jep.35.4.197","locator":"Conclusion, printed pp. 216–217; introduction p. 200","evidence_summary":"The authors explicitly identify mechanism weights, general-equilibrium scaling of mobility programs, and effects of revitalization on preexisting residents as open issues."},"current_reference":{"authors":["Lawrence F. Katz"],"title":"Neighborhood Effects and Missing Markets for Opportunity","year":2026,"url":"https://www.nber.org/papers/w35419","doi":"10.3386/w35419","locator":"Abstract","evidence_summary":"The abstract discusses missing markets and mobility assistance; the older Chyn–Katz review explicitly establishes the scaling and mechanism research agendas."},"formalization_note":"The source explicitly identifies the underlying gap or agenda. Population, horizon, interventions, estimands, and auxiliary assumptions are proposed operational choices, not a canonical source theorem. Partial later answers are compatible with this scoped question; the citation alone does not prove absence of every subsequent solution. The 2021 Chyn–Katz review supplies an earlier verified agenda than the original catalogue’s 2023 or 2026 supporting citation."},"statusQualification":"Published question or research agenda verified at the cited locator. The mathematical specification is adapted for this catalogue and includes additional assumptions; source publication does not establish first-ever priority or certify this exact model remains unsolved. The source explicitly identifies the underlying gap or agenda. Population, horizon, interventions, estimands, and auxiliary assumptions are proposed operational choices, not a canonical source theorem. Partial later answers are compatible with this scoped question; the citation alone does not prove absence of every subsequent solution. The 2021 Chyn–Katz review supplies an earlier verified agenda than the original catalogue’s 2023 or 2026 supporting citation.","statusReviewedAt":"2026-10-08"}
% FIRST_POSTED: 2026-10-08
% ENTRY_KIND: statement_only
% !TeX program = lualatex
\documentclass[11pt]{article}
\usepackage{fontspec}
\usepackage[a4paper,margin=25mm]{geometry}
\usepackage{amsmath,amssymb,mathtools}
\usepackage{xurl}
\usepackage[hidelinks]{hyperref}
\newcommand{\sref}[2]{\hyperlink{source:#1:#2}{[S#2]}}
\setlength{\emergencystretch}{3em}
\begin{document}
% problemId:30007746
\section*{Scaling housing mobility programs}\label{30007746}
\subsection{Definitions and mathematical statement}
Consider low-income families eligible for housing vouchers in one metropolitan area. Let \(s\in[0,1]\) be the share of eligible families offered an assistance package combining search support and moving subsidies. Housing prices, available units, landlord acceptance, neighborhood composition, and school capacity may respond to \(s\). For baseline family \(i\), define \(Y_i(s)\) as children’s discounted adult earnings plus a stated monetary-equivalent measure of family health and housing utility; define \(K(s)\) program resource cost. For baseline noneligible household \(j\), let \(V_j(s)\) be monetary-equivalent utility. The scaling target is
\[W(s)=E\!\left[\sum_iY_i(s)+\sum_jV_j(s)\right]-K(s).\]
Estimate \(W(s)-W(0)\) at prespecified saturation levels rather than multiply a small trial’s average effect by the eligible population. Use randomized market saturation where feasible, or a spatial model disciplined by experimental moves and identified supply responses. Specify school and peer spillovers, moving disruption, and follow-up extrapolation. Report distributional effects and uncertainty in general-equilibrium parameters. The Chyn–Katz conclusion explicitly identifies scaling mobility programs as an unresolved research agenda; the saturation design and welfare definition operationalize that agenda.

\par\textbf{Formulation and scope.} The source explicitly identifies the underlying gap or agenda. Population, horizon, interventions, estimands, and auxiliary assumptions are proposed operational choices, not a canonical source theorem. Partial later answers are compatible with this scoped question; the citation alone does not prove absence of every subsequent solution. The 2021 Chyn–Katz review supplies an earlier verified agenda than the original catalogue’s 2023 or 2026 supporting citation.

\par\textbf{Open-question provenance.} An explicit question or research agenda was verified at the source locator. This does not by itself certify that every later version remains unresolved \sref{30007746}{1}.

\par\textbf{Historical priority.} Earliest explicit statement located in this audit; no exhaustive priority search or first-ever claim. The mathematical specification below is adapted, not quoted from the source.

\subsection{Short English statement}
Estimate welfare at different saturation levels of a housing-mobility assistance program after rents, neighborhood composition, and school capacity respond.

\subsection{Sources}
\begin{itemize}\raggedright
\item[\textnormal{[S1]}]\hypertarget{source:30007746:1}{} Eric Chyn, Lawrence F. Katz. 2021. Neighborhoods Matter: Assessing the Evidence for Place Effects. Conclusion, printed pp. 216–217; introduction p. 200.
\url{https://lkatz.scholars.harvard.edu/sites/g/files/omnuum5961/files/lkatz/files/chyn_katz_jep_2021_all.pdf}.
DOI: 10.1257/jep.35.4.197.
{\small\textit{Earliest verified question/agenda reference; historical priority qualified below. The authors explicitly identify mechanism weights, general-equilibrium scaling of mobility programs, and effects of revitalization on preexisting residents as open issues.}}

\item[\textnormal{[S2]}]\hypertarget{source:30007746:2}{} Lawrence F. Katz. 2026. Neighborhood Effects and Missing Markets for Opportunity. Abstract.
\url{https://www.nber.org/papers/w35419}.
DOI: 10.3386/w35419.
{\small\textit{Supporting or current literature; see evidence qualification. The abstract discusses missing markets and mobility assistance; the older Chyn–Katz review explicitly establishes the scaling and mechanism research agendas.}}
\end{itemize}
\end{document}
